\documentclass[../report_main.tex]{subfiles}
\begin{document}

\section{Summary of Achievements}
This project successfully culminated in the development of a sophisticated, microscopic smart traffic simulation framework written entirely in the Java programming language. The resulting application effectively models the complex interplay of vehicle kinematics, localized driver decision-making, and geometric collision detection within a highly customizable urban map environment. By meticulously implementing fundamental physics equations for motion and leveraging mathematical constructs such as Cramer's rule for precise spatial calculations, the simulation provides an incredibly realistic and computationally efficient representation of urban traffic flow. The strategic inclusion of diverse vehicle types, ranging from standard civilian cars to high-priority emergency vehicles, alongside distinct behavioral profiles, adds significant depth and realism to the simulation engine. Ultimately, this structural robustness and visual clarity establish the application as a highly viable educational tool for understanding traffic dynamics and software engineering paradigms.

\section{Future Work}
While the current implementation successfully fulfills and exceeds the primary project objectives, several advanced enhancements could be explored in future development iterations to further expand the capabilities of the simulation. A primary area for expansion involves the integration of machine learning techniques, specifically implementing reinforcement learning algorithms to dynamically train autonomous vehicle agents within the safe confines of the simulation environment. Another significant enhancement would be the development of an advanced traffic management module capable of dynamically scheduling traffic light phases based on real-time traffic density metrics gathered from simulated sensors, rather than relying on static timers. Furthermore, the introduction of pedestrian agents and interactive crosswalks would vastly improve the realism of the urban model by simulating the intricate and often unpredictable interactions between vehicular traffic and foot traffic. Finally, incorporating comprehensive data analytics and export features would allow researchers to systematically analyze aggregate metrics such as average intersection wait times, total road throughput, and the spontaneous formation of congestion hotspots over extended simulation periods.

\end{document}
